\section{Przywrócenie wszystkich punktów sfery}
Pozostało uzasadnić $S^2\equiv X$. Nie wystarcza zdanie „zbiór przeliczalny jest mały”. Potrzebujemy skończonego rozkładu i obrotów przypisanych jego częściom.

\subsection{Wybór kąta poza zbiorem przeliczalnym}
\begin{lemat}\label{lem:nieprzeliczalnosc}
Żaden niepusty otwarty przedział liczb rzeczywistych nie jest przeliczalny. W szczególności można w nim wybrać liczbę poza dowolnym podanym zbiorem co najwyżej przeliczalnym.
\end{lemat}
\begin{proof}
Niech $x_1,x_2,\ldots$ będzie dowolną listą punktów przedziału. Wybierzmy w nim niejednopunktowy przedział domknięty $I_0$. Indukcyjnie wybieramy niejednopunktowy przedział domknięty $I_n\subset I_{n-1}$ o długości co najwyżej połowy poprzedniej, taki że $x_n\notin I_n$. Możemy go umieścić po jednej stronie $x_n$; jeśli $x_n$ nie należy do $I_{n-1}$, dowolny mniejszy podprzedział wystarczy.

Przedziały są zstępujące, domknięte i ich długości dążą do zera. Z zupełności liczb rzeczywistych ich przecięcie zawiera dokładnie jeden punkt $x$: jest nim wspólna granica lewych i prawych końców. Dla każdego $n$ mamy $x\in I_n$, więc $x\ne x_n$. Lista nie obejmowała całego przedziału.
\end{proof}
Sfera też jest nieprzeliczalna: zawiera półokrąg parametryzowany różnowartościowo przez $\theta\mapsto(\cos\theta,\sin\theta,0)$ dla $0<\theta<\pi$. Przeliczalność zbioru trójek indeksów naturalnych wynika z wyliczania ich według sumy indeksów; dla każdej sumy jest tylko skończenie wiele trójek. Dlatego suma przeliczalnej rodziny skończonych zbiorów wykluczonych kątów jest przeliczalna.

\subsection{Obrót rozdzielający przeliczalny zbiór}
\begin{lemat}\label{lem:obrot-rozdzielajacy}
Dla dowolnego co najwyżej przeliczalnego $D\subseteq S^2$ istnieje obrót $R\in\mathrm{SO}(3)$, dla którego zbiory $D,RD,R^2D,\ldots$ są parami rozłączne.
\end{lemat}
\begin{proof}
Jeśli $D=\varnothing$, teza jest oczywista. Niech $D\ne\varnothing$. Wybieramy punkt $u\in S^2$ poza $D\cup(-D)$, co jest możliwe, bo ten ostatni zbiór jest przeliczalny. Żaden punkt $D$ nie leży na osi $\R u$.

Niech $R_u(\theta)$ oznacza obrót wokół tej osi o kąt $\theta$. Można go zdefiniować, wybierając ortonormalną bazę $u,v,w$ o dodatniej orientacji i przyjmując
\[
 R_u(\theta)(\alpha u+\beta v+\gamma w)=\alpha u+
 (\beta\cos\theta-\gamma\sin\theta)v+
 (\beta\sin\theta+\gamma\cos\theta)w.
\]
Ze wzorów dodawania kątów wynika $R_u(\theta)^n=R_u(n\theta)$.

Ustalmy $d,e\in D$ i $n\ge1$. Ponieważ $d$ nie leży na osi, jego składowa w płaszczyźnie $u^\perp$ jest niezerowa. Obracając tę składową wokół osi, otrzymujemy okrąg; każde położenie na nim wyznacza kąt jednoznacznie modulo $2\pi$. Zatem albo nie istnieje obrót wokół tej osi przenoszący $d$ na $e$, albo istnieje dokładnie jedna klasa kątów $\alpha$ modulo $2\pi$ o tej własności.

Równanie $R_u(n\theta)d=e$ sprowadza się w drugim przypadku do $n\theta\equiv\alpha\pmod{2\pi}$. W przedziale $0<\theta<2\pi$ ma najwyżej $n$ rozwiązań. Dla każdej trójki $(d,e,n)$ wykluczamy więc skończenie wiele kątów. Trójek jest co najwyżej przeliczalnie wiele. Na mocy lematu~\ref{lem:nieprzeliczalnosc} możemy wybrać $\theta$ poza unią wszystkich wykluczonych kątów.

Dla $R=R_u(\theta)$ mamy teraz
\begin{equation}\label{eq:7-1}
 R^nD\cap D=\varnothing\qquad(n\ge1).
\end{equation}
Jeśli dla $0\le i<j$ istniałby punkt w $R^iD\cap R^jD$, to jego obraz przez $R^{-i}$ należałby do $D\cap R^{j-i}D$, wbrew \eqref{eq:7-1}. Dostajemy żądaną rozłączność wszystkich poziomów.
\end{proof}
\begin{uwaga}
Nowy obrót $R$ nie musi należeć do $G$. Ponieważ $D$ jest $G$-niezmienniczy, żaden element $G$ nie mógłby przenieść niepustego $D$ całkowicie poza $D$. Od tego miejsca dla równoważności całej sfery dopuszczamy pełną grupę obrotów $\mathrm{SO}(3)$.
\end{uwaga}
\begin{twierdzenie}[Pełna sfera jest paradoksalna]\label{tw:pelna-sfera}
Sfera $S^2$ jest równoważna przez skończony rozkład dwóm swoim etykietowanym kopiom, przy użyciu samych obrotów wokół $0$. Wystarcza co najwyżej $16$ części.
\end{twierdzenie}
\begin{proof}
Stosujemy lemat~\ref{lem:obrot-rozdzielajacy} do zbioru $D$ z \eqref{eq:6-1}. Niech $L=\bigcup_{n=0}^\infty R^nD$. Wszystkie poziomy leżą na sferze, więc lemat~\ref{lem:lancuch} daje dwuczęściowy rozkład $S^2\equiv X$. Twierdzenie~\ref{tw:paradoksalnosc-x} daje $X\equiv2X$. Przeniesienie paradoksalności z wniosku~\ref{wn:przenoszenie} kończy dowód istnienia rozkładu $S^2\equiv2S^2$.

Poniżej podajemy dokładne części, aby samo odwołanie do przechodniości nie ukrywało konstrukcji ani oszacowania ich liczby. Niech $E_0=L$, $E_1=S^2\setminus L$, $g_0=R$, $g_1=I$. Bijekcja $f:S^2\to X$ ma na $E_i$ postać $g_i$. Zbiory $g_0E_0=L\setminus D$ i $g_1E_1=S^2\setminus L$ tworzą partycję $X$.

Oznaczmy przez $C_j,h_j,\ell(j)$ cztery części, ruchy i etykiety z dowodu twierdzenia~\ref{tw:paradoksalnosc-x}. Dla $i,k\in\{0,1\}$ i $j\in\{1,2,3,4\}$ połóżmy
\begin{equation}\label{eq:7-2}
 H_{ijk}=\{x\in E_i:g_ix\in C_j,\ h_jg_ix\in g_kE_k\}.
\end{equation}
Na tej części działa obrót
\begin{equation}\label{eq:7-3}
 q_{ijk}=g_k^{-1}h_jg_i,
\end{equation}
z etykietą docelową $\ell(j)$. Dla każdego $x$ indeksy $i,j,k$ są wyznaczone jednoznacznie, w tej kolejności. Zatem części \eqref{eq:7-2} tworzą partycję sfery.

Sprawdźmy jeszcze obrazy. Całe odwzorowanie jest złożeniem bijekcji $f$, czteroczęściowej bijekcji $X\to2X$ i odwrotnej bijekcji $f^{-1}$ w każdej kopii. Jest zatem bijekcją $S^2\to2S^2$. Stąd obrazy kawałków są rozłączne wewnątrz każdej etykiety i pokrywają obie sfery. Indeksów jest $2\cdot4\cdot2=16$. Każde $q_{ijk}$ jest złożeniem obrotów, a więc nadal jednym obrotem w $\mathrm{SO}(3)$.
\end{proof}

